
\documentclass{article}
\usepackage{colortbl}
\usepackage{makecell}
\usepackage{multirow}
\usepackage{supertabular}

\begin{document}

\newcounter{utterance}

\centering \large Interaction Transcript for game `cladder', experiment `full\_v1.5\_default', episode 53 with Qwen3.5{-}27B.
\vspace{24pt}

{ \footnotesize  \setcounter{utterance}{1}
\setlength{\tabcolsep}{0pt}
\begin{supertabular}{c@{$\;$}|p{.15\linewidth}@{}p{.15\linewidth}p{.15\linewidth}p{.15\linewidth}p{.15\linewidth}p{.15\linewidth}}
    \# & \multicolumn{2}{c}{Player} && \multicolumn{2}{c}{Game Master} \\
    \hline

    \theutterance \stepcounter{utterance}  
    & & & \multicolumn{4}{p{0.6\linewidth}}{
        \cellcolor[rgb]{0.9,0.9,0.9}{
            \makecell[{{p{\linewidth}}}]{
                \texttt{\tiny{[P1$\langle$GM]}}
                \texttt{You are an expert in causal inference. The following question is not a typical commonsense query, but rather a meticulously designed question created by a professor specializing in causal inference, intended to assess the students' mastery of the course content.} \\
\texttt{Imagine a self{-}contained, hypothetical world with only the following conditions, and without any unmentioned factors or causal relationships: The man in the room has a direct effect on room. The candle has a direct effect on room. Method 1: We look directly at how the man in the room correlates with room in general. Method 2: We look at this correlation case by case according to the candle. To understand how the man in the room affects room, is it more correct to use the Method 1 than Method 2?} \\
\texttt{Start your answer with 'Yes' or 'No', followed by additional reasoning or evidence to support your explanation.} \\
            }
        }
    }
     \\ \\

    \theutterance \stepcounter{utterance}  
    & \multicolumn{4}{p{0.6\linewidth}}{
        \cellcolor[rgb]{0.561,0.737,0.561}{
            \makecell[{{p{\linewidth}}}]{
                \texttt{\tiny{[P1$\rangle$GM]}}
                \texttt{No} \\
\\ 
\texttt{In this hypothetical scenario, the causal structure is defined as two independent causes (the man and the candle) both directly affecting a single outcome (the room). This structure represents a "common effect" or "collider" configuration where the room is the collider, but more importantly for the estimation of the man's effect, the man and the candle are not causally related to each other, nor is the candle a confounder (a common cause of both the man and the room).} \\
\\ 
\texttt{Method 1 calculates the marginal association between the man and the room. Since there are no unmentioned confounders and the candle is not a confounder, this marginal association provides an unbiased estimate of the total causal effect of the man on the room.} \\
\\ 
\texttt{Method 2 involves conditioning on the candle (stratifying by the candle's state). While the candle is a cause of the room, it is not a confounder of the relationship between the man and the room. Conditioning on a variable that is a cause of the outcome but not a confounder (often called a "post{-}treatment" variable in some contexts, or simply a non{-}confounder here) does not introduce bias in this specific DAG structure, but it is statistically inefficient. It reduces the sample size within each stratum, increasing the variance of the estimate without reducing bias. Furthermore, if the candle and the man were correlated (which is not stated, but possible in observational data), conditioning on the candle could introduce bias if the candle were a collider, but here the candle is a direct cause. However, the most critical point in causal inference pedagogy for this specific setup is that **conditioning is only necessary to block backdoor paths from confounders**. Since there are no backdoor paths from the man to the room (the only path is the direct causal arrow), conditioning on the candle is unnecessary.} \\
\\ 
\texttt{Therefore, Method 1 is "more correct" in the sense that it is the simpler, more efficient, and sufficient approach to identify the causal effect without introducing unnecessary complexity or variance. Method 2 is not "wrong" in terms of bias (assuming no interaction effects that require stratification for interpretation), but it is not *more* correct than Method 1; in fact, Method 1 is the standard approach for estimating the total effect when no confounding exists. If the question implies whether Method 2 is *required* or *superior* to Method 1, the answer is definitively no. Method 1 is the correct and preferred method here.} \\
            }
        }
    }
    & & \\ \\

    \theutterance \stepcounter{utterance}  
    & & & \multicolumn{2}{p{0.3\linewidth}}{
        \cellcolor[rgb]{0.9,0.9,0.9}{
            \makecell[{{p{\linewidth}}}]{
                \texttt{\tiny{[GM$|$GM]}}
                \texttt{No} \\
            }
        }
    }
    & & \\ \\

    \theutterance \stepcounter{utterance}  
    & & & \multicolumn{2}{p{0.3\linewidth}}{
        \cellcolor[rgb]{0.9,0.9,0.9}{
            \makecell[{{p{\linewidth}}}]{
                \texttt{\tiny{[GM$|$GM]}}
                \texttt{yes} \\
            }
        }
    }
    & & \\ \\

    \theutterance \stepcounter{utterance}  
    & & & \multicolumn{2}{p{0.3\linewidth}}{
        \cellcolor[rgb]{0.9,0.9,0.9}{
            \makecell[{{p{\linewidth}}}]{
                \texttt{\tiny{[GM$|$GM]}}
                \texttt{game\_result = LOSE} \\
            }
        }
    }
    & & \\ \\

\end{supertabular}
}

\end{document}
